Financial privacy evaluation must specify what to detect as well as how it appears in text. Korean financial documents---including call-center transcripts and scanned forms---combine jurisdiction-specific categories, noncanonical identifiers and contextual attributes. Recognizing a standard identifier pattern therefore addresses only part of the task. We introduce \textbf{Kiii-Kiii} (dataset \kiii{}), a benchmark that makes these dimensions explicit. A 36-category taxonomy distinguishes 24 categories with statutory grounds from 12 additional quasi-identifier and profiling categories, without equating detection with a legal duty to redact. The synthetic preview contains 1,440 documents and 218,664 spans across ten document types, with code-generated identifiers, applicable checksum checks and four levels of surface-form variation. A matched full-context versus local-window protocol holds target regions and output budgets fixed, enabling analysis of whether additional context improves span recovery. Exact span scoring retains malformed and truncated responses as failures. Across 12 language models and three frozen baselines, strict extraction is limited by both output-contract failures and missed spans. A fixed, post-hoc itemwise parser provides a closer observable proxy for detection, reaching 31.58 F1 while retaining all gold and penalizing invalid items. Local context improves this diagnostic in all ten completed paired comparisons; strict scores favor local in nine. These findings motivate separate reporting of extraction reliability and detectable span recovery. The preview is released under CC BY-NC 4.0, with a formal release to follow.
